% 5.1 抽象群
\section{抽象群}

对空间群布里渊区中的任一波矢 $\mathbf{k}$，小群 $\mathbf{G}^{\mathbf{k}}$ 是一个无限群。然而，当 $\mathbf{k}$ 为对称点时，这个无限群可以联系于一个有限群 ${}^{H}\mathbf{G}^{\mathbf{k}} = \mathbf{G}^{\mathbf{k}}/\mathbf{T}^{\mathbf{k}}$（见 3.8 节）；其他情形则联系于小陪群的中央延拓群 $\overline{\mathbf{G}}^{\mathbf{k}*}$（见定理 3.7.2）。人们常常发现，有若干个群 ${}^{H}\mathbf{G}^{\mathbf{k}}$ 或 $\overline{\mathbf{G}}^{\mathbf{k}*}$——也许属于同一空间群布里渊区中的不同波矢，或属于不同空间群布里渊区中的同一波矢，或属于不同空间群布里渊区中的不同波矢——全都同构于同一个抽象群。由于同一抽象群在许多不同空间群中反复出现，我们在表 5.1 中完整地列出了（空间群的（单值与双值）表示中出现的）所有抽象群的全部不可约表示。

每个抽象群由一组生成元 $P, Q, R, S, \dots$ 刻画，群的每个元素都可以写成 $P^{\alpha}Q^{\beta}R^{\gamma}S^{\delta}\dots$ 的形式。只要给出群的乘法表，群就完全确定；而为了确定乘法表，只需有元素 $P, Q, R, S$ 等的定义关系或生成关系即可（Coxeter and Moser 1965, Hall and Senior 1964, Miller 1894, 1946）。例 1.2.1 对一个非常简单的群演示过这一点。在表 5.1 中我们给出了生成关系、元素在群的各类中的划分以及群的特征标表。每个抽象群的表示标记为 $R_{1}, R_{2}, \dots, R_{r}$，并且对每个简并表示我们给出其生成元的矩阵代表。任何给定简并表示的全套矩阵代表随之可以确定。

\begin{figure}[H]
\centering
\includegraphics[width=0.86\textwidth]{c5_t51_p01.jpg}
\caption*{\textbf{表 5.1}\ 抽象群的定义关系、类、特征标表与矩阵代表（230 个空间群的表示表所需；原书扫描，第 1 页，共 60 页）}
\end{figure}
\begin{figure}[H]
\centering
\includegraphics[width=0.86\textwidth]{c5_t51_p02.jpg}
\caption*{\textbf{表 5.1}（第 2 页，共 60 页；原书第 227 页）}
\end{figure}
\begin{figure}[H]
\centering
\includegraphics[width=0.86\textwidth]{c5_t51_p03.jpg}
\caption*{\textbf{表 5.1}（第 3 页，共 60 页；原书第 228 页）}
\end{figure}
\begin{figure}[H]
\centering
\includegraphics[width=0.86\textwidth]{c5_t51_p04.jpg}
\caption*{\textbf{表 5.1}（第 4 页，共 60 页；原书第 229 页）}
\end{figure}
\begin{figure}[H]
\centering
\includegraphics[width=0.86\textwidth]{c5_t51_p05.jpg}
\caption*{\textbf{表 5.1}（第 5 页，共 60 页；原书第 230 页）}
\end{figure}
\begin{figure}[H]
\centering
\includegraphics[width=0.86\textwidth]{c5_t51_p06.jpg}
\caption*{\textbf{表 5.1}（第 6 页，共 60 页；原书第 231 页）}
\end{figure}
\begin{figure}[H]
\centering
\includegraphics[width=0.86\textwidth]{c5_t51_p07.jpg}
\caption*{\textbf{表 5.1}（第 7 页，共 60 页；原书第 232 页）}
\end{figure}
\begin{figure}[H]
\centering
\includegraphics[width=0.86\textwidth]{c5_t51_p08.jpg}
\caption*{\textbf{表 5.1}（第 8 页，共 60 页；原书第 233 页）}
\end{figure}
\begin{figure}[H]
\centering
\includegraphics[width=0.86\textwidth]{c5_t51_p09.jpg}
\caption*{\textbf{表 5.1}（第 9 页，共 60 页；原书第 234 页）}
\end{figure}
\begin{figure}[H]
\centering
\includegraphics[width=0.86\textwidth]{c5_t51_p10.jpg}
\caption*{\textbf{表 5.1}（第 10 页，共 60 页；原书第 235 页）}
\end{figure}
\begin{figure}[H]
\centering
\includegraphics[width=0.86\textwidth]{c5_t51_p11.jpg}
\caption*{\textbf{表 5.1}（第 11 页，共 60 页；原书第 236 页）}
\end{figure}
\begin{figure}[H]
\centering
\includegraphics[width=0.86\textwidth]{c5_t51_p12.jpg}
\caption*{\textbf{表 5.1}（第 12 页，共 60 页；原书第 237 页）}
\end{figure}
\begin{figure}[H]
\centering
\includegraphics[width=0.86\textwidth]{c5_t51_p13.jpg}
\caption*{\textbf{表 5.1}（第 13 页，共 60 页；原书第 238 页）}
\end{figure}
\begin{figure}[H]
\centering
\includegraphics[width=0.86\textwidth]{c5_t51_p14.jpg}
\caption*{\textbf{表 5.1}（第 14 页，共 60 页；原书第 239 页）}
\end{figure}
\begin{figure}[H]
\centering
\includegraphics[width=0.86\textwidth]{c5_t51_p15.jpg}
\caption*{\textbf{表 5.1}（第 15 页，共 60 页；原书第 240 页）}
\end{figure}
\begin{figure}[H]
\centering
\includegraphics[width=0.86\textwidth]{c5_t51_p16.jpg}
\caption*{\textbf{表 5.1}（第 16 页，共 60 页；原书第 241 页）}
\end{figure}
\begin{figure}[H]
\centering
\includegraphics[width=0.86\textwidth]{c5_t51_p17.jpg}
\caption*{\textbf{表 5.1}（第 17 页，共 60 页；原书第 242 页）}
\end{figure}
\begin{figure}[H]
\centering
\includegraphics[width=0.86\textwidth]{c5_t51_p18.jpg}
\caption*{\textbf{表 5.1}（第 18 页，共 60 页；原书第 243 页）}
\end{figure}
\begin{figure}[H]
\centering
\includegraphics[width=0.86\textwidth]{c5_t51_p19.jpg}
\caption*{\textbf{表 5.1}（第 19 页，共 60 页；原书第 244 页）}
\end{figure}
\begin{figure}[H]
\centering
\includegraphics[width=0.86\textwidth]{c5_t51_p20.jpg}
\caption*{\textbf{表 5.1}（第 20 页，共 60 页；原书第 245 页）}
\end{figure}
\begin{figure}[H]
\centering
\includegraphics[width=0.86\textwidth]{c5_t51_p21.jpg}
\caption*{\textbf{表 5.1}（第 21 页，共 60 页；原书第 246 页）}
\end{figure}
\begin{figure}[H]
\centering
\includegraphics[width=0.86\textwidth]{c5_t51_p22.jpg}
\caption*{\textbf{表 5.1}（第 22 页，共 60 页；原书第 247 页）}
\end{figure}
\begin{figure}[H]
\centering
\includegraphics[width=0.86\textwidth]{c5_t51_p23.jpg}
\caption*{\textbf{表 5.1}（第 23 页，共 60 页；原书第 248 页）}
\end{figure}
\begin{figure}[H]
\centering
\includegraphics[width=0.86\textwidth]{c5_t51_p24.jpg}
\caption*{\textbf{表 5.1}（第 24 页，共 60 页；原书第 249 页）}
\end{figure}
\begin{figure}[H]
\centering
\includegraphics[width=0.86\textwidth]{c5_t51_p25.jpg}
\caption*{\textbf{表 5.1}（第 25 页，共 60 页；原书第 250 页）}
\end{figure}
\begin{figure}[H]
\centering
\includegraphics[width=0.86\textwidth]{c5_t51_p26.jpg}
\caption*{\textbf{表 5.1}（第 26 页，共 60 页；原书第 251 页）}
\end{figure}
\begin{figure}[H]
\centering
\includegraphics[width=0.86\textwidth]{c5_t51_p27.jpg}
\caption*{\textbf{表 5.1}（第 27 页，共 60 页；原书第 252 页）}
\end{figure}
\begin{figure}[H]
\centering
\includegraphics[width=0.86\textwidth]{c5_t51_p28.jpg}
\caption*{\textbf{表 5.1}（第 28 页，共 60 页；原书第 253 页）}
\end{figure}
\begin{figure}[H]
\centering
\includegraphics[width=0.86\textwidth]{c5_t51_p29.jpg}
\caption*{\textbf{表 5.1}（第 29 页，共 60 页；原书第 254 页）}
\end{figure}
\begin{figure}[H]
\centering
\includegraphics[width=0.86\textwidth]{c5_t51_p30.jpg}
\caption*{\textbf{表 5.1}（第 30 页，共 60 页；原书第 255 页）}
\end{figure}
\begin{figure}[H]
\centering
\includegraphics[width=0.86\textwidth]{c5_t51_p31.jpg}
\caption*{\textbf{表 5.1}（第 31 页，共 60 页；原书第 256 页）}
\end{figure}
\begin{figure}[H]
\centering
\includegraphics[width=0.86\textwidth]{c5_t51_p32.jpg}
\caption*{\textbf{表 5.1}（第 32 页，共 60 页；原书第 257 页）}
\end{figure}
\begin{figure}[H]
\centering
\includegraphics[width=0.86\textwidth]{c5_t51_p33.jpg}
\caption*{\textbf{表 5.1}（第 33 页，共 60 页；原书第 258 页）}
\end{figure}
\begin{figure}[H]
\centering
\includegraphics[width=0.86\textwidth]{c5_t51_p34.jpg}
\caption*{\textbf{表 5.1}（第 34 页，共 60 页；原书第 259 页）}
\end{figure}
\begin{figure}[H]
\centering
\includegraphics[width=0.86\textwidth]{c5_t51_p35.jpg}
\caption*{\textbf{表 5.1}（第 35 页，共 60 页；原书第 260 页）}
\end{figure}
\begin{figure}[H]
\centering
\includegraphics[width=0.86\textwidth]{c5_t51_p36.jpg}
\caption*{\textbf{表 5.1}（第 36 页，共 60 页；原书第 261 页）}
\end{figure}
\begin{figure}[H]
\centering
\includegraphics[width=0.86\textwidth]{c5_t51_p37.jpg}
\caption*{\textbf{表 5.1}（第 37 页，共 60 页；原书第 262 页）}
\end{figure}
\begin{figure}[H]
\centering
\includegraphics[width=0.86\textwidth]{c5_t51_p38.jpg}
\caption*{\textbf{表 5.1}（第 38 页，共 60 页；原书第 263 页）}
\end{figure}
\begin{figure}[H]
\centering
\includegraphics[width=0.86\textwidth]{c5_t51_p39.jpg}
\caption*{\textbf{表 5.1}（第 39 页，共 60 页；原书第 264 页）}
\end{figure}
\begin{figure}[H]
\centering
\includegraphics[width=0.86\textwidth]{c5_t51_p40.jpg}
\caption*{\textbf{表 5.1}（第 40 页，共 60 页；原书第 265 页）}
\end{figure}
\begin{figure}[H]
\centering
\includegraphics[width=0.86\textwidth]{c5_t51_p41.jpg}
\caption*{\textbf{表 5.1}（第 41 页，共 60 页；原书第 266 页）}
\end{figure}
\begin{figure}[H]
\centering
\includegraphics[width=0.86\textwidth]{c5_t51_p42.jpg}
\caption*{\textbf{表 5.1}（第 42 页，共 60 页；原书第 267 页）}
\end{figure}
\begin{figure}[H]
\centering
\includegraphics[width=0.86\textwidth]{c5_t51_p43.jpg}
\caption*{\textbf{表 5.1}（第 43 页，共 60 页；原书第 268 页）}
\end{figure}
\begin{figure}[H]
\centering
\includegraphics[width=0.86\textwidth]{c5_t51_p44.jpg}
\caption*{\textbf{表 5.1}（第 44 页，共 60 页；原书第 269 页）}
\end{figure}
\begin{figure}[H]
\centering
\includegraphics[width=0.86\textwidth]{c5_t51_p45.jpg}
\caption*{\textbf{表 5.1}（第 45 页，共 60 页；原书第 270 页）}
\end{figure}
\begin{figure}[H]
\centering
\includegraphics[width=0.86\textwidth]{c5_t51_p46.jpg}
\caption*{\textbf{表 5.1}（第 46 页，共 60 页；原书第 271 页）}
\end{figure}
\begin{figure}[H]
\centering
\includegraphics[width=0.86\textwidth]{c5_t51_p47.jpg}
\caption*{\textbf{表 5.1}（第 47 页，共 60 页；原书第 272 页）}
\end{figure}
\begin{figure}[H]
\centering
\includegraphics[width=0.86\textwidth]{c5_t51_p48.jpg}
\caption*{\textbf{表 5.1}（第 48 页，共 60 页；原书第 273 页）}
\end{figure}
\begin{figure}[H]
\centering
\includegraphics[width=0.86\textwidth]{c5_t51_p49.jpg}
\caption*{\textbf{表 5.1}（第 49 页，共 60 页；原书第 274 页）}
\end{figure}
\begin{figure}[H]
\centering
\includegraphics[width=0.86\textwidth]{c5_t51_p50.jpg}
\caption*{\textbf{表 5.1}（第 50 页，共 60 页；原书第 275 页）}
\end{figure}
\begin{figure}[H]
\centering
\includegraphics[width=0.86\textwidth]{c5_t51_p51.jpg}
\caption*{\textbf{表 5.1}（第 51 页，共 60 页；原书第 276 页）}
\end{figure}
\begin{figure}[H]
\centering
\includegraphics[width=0.86\textwidth]{c5_t51_p52.jpg}
\caption*{\textbf{表 5.1}（第 52 页，共 60 页；原书第 277 页）}
\end{figure}
\begin{figure}[H]
\centering
\includegraphics[width=0.86\textwidth]{c5_t51_p53.jpg}
\caption*{\textbf{表 5.1}（第 53 页，共 60 页；原书第 278 页）}
\end{figure}
\begin{figure}[H]
\centering
\includegraphics[width=0.86\textwidth]{c5_t51_p54.jpg}
\caption*{\textbf{表 5.1}（第 54 页，共 60 页；原书第 279 页）}
\end{figure}
\begin{figure}[H]
\centering
\includegraphics[width=0.86\textwidth]{c5_t51_p55.jpg}
\caption*{\textbf{表 5.1}（第 55 页，共 60 页；原书第 280 页）}
\end{figure}
\begin{figure}[H]
\centering
\includegraphics[width=0.86\textwidth]{c5_t51_p56.jpg}
\caption*{\textbf{表 5.1}（第 56 页，共 60 页；原书第 281 页）}
\end{figure}
\begin{figure}[H]
\centering
\includegraphics[width=0.86\textwidth]{c5_t51_p57.jpg}
\caption*{\textbf{表 5.1}（第 57 页，共 60 页；原书第 282 页）}
\end{figure}
\begin{figure}[H]
\centering
\includegraphics[width=0.86\textwidth]{c5_t51_p58.jpg}
\caption*{\textbf{表 5.1}（第 58 页，共 60 页；原书第 283 页）}
\end{figure}
\begin{figure}[H]
\centering
\includegraphics[width=0.86\textwidth]{c5_t51_p59.jpg}
\caption*{\textbf{表 5.1}（第 59 页，共 60 页；原书第 284 页）}
\end{figure}
\begin{figure}[H]
\centering
\includegraphics[width=0.86\textwidth]{c5_t51_p60.jpg}
\caption*{\textbf{表 5.1}（第 60 页，共 60 页；原书第 285 页）}
\end{figure}

可以把（其特征标表已在表 2.2 中给出的）每个晶体学点群与这些抽象群之一对应起来（Belova, Belov, and Shubnikov 1948）。例如，抽象群 $\mathbf{G}^{4}_{8}$ 同构于点群 $422\ (D_{4})$。通过观察不难看出点群 $422\ (D_{4})$ 的哪些元素对应于 $\mathbf{G}^{4}_{8}$ 的生成元，特别是因为 $P$ 与 $Q$ 的选取受到表示 $E$ 的矩阵的限制——由表 5.1 的钥匙（key），
\begin{equation*}
P = \begin{pmatrix} 0 & 1\\ -1 & 0 \end{pmatrix}, \quad Q = \begin{pmatrix} 1 & 0\\ 0 & -1 \end{pmatrix}.
\end{equation*}
因此由表 2.3，$P$ 必是 $C^{-}_{4z}$，$Q$ 必是 $C_{2x}$，于是两个群之间的完整对应为
\begin{center}
\begin{tabular}{lcc}
 & $422\ (D_{4})$ & $\mathbf{G}^{4}_{8}$\\
$C_{1}$ & $E$ & $E$\\
$C_{2}$ & $C_{2z}$ & $P^{2}$\\
$C_{3}$ & $\begin{array}{c}C^{-}_{4z}\\ C^{+}_{4z}\end{array}$ & $\begin{array}{c}P\\ P^{3}\end{array}$\\
$C_{4}$ & $\begin{array}{c}C_{2x}\\ C_{2y}\end{array}$ & $\begin{array}{c}Q\\ P^{2}Q\end{array}$\\
$C_{5}$ & $\begin{array}{c}C_{2b}\\ C_{2a}\end{array}$ & $\begin{array}{c}PQ\\ P^{3}Q\end{array}$
\end{tabular}
\end{center}
当然，元素的认定必须恰当顾及抽象群与其作为点群的实现之间的同构，例如 $PQ = C^{-}_{4z}C_{2x} = C_{2b}$（用 Jones 符号或表 1.5）。

每个晶体学点群作为表 5.1 中某一抽象群的一种体现的认定由表 5.2 给出；这一认定并不总是唯一的，且某些情形下矩阵代表与第二章所给者相差两个群之间的一个幺正变换。
